\documentclass[11pt,a4paper]{article}
\usepackage[utf8]{inputenc}
\usepackage[T1]{fontenc}
\usepackage[spanish,es-nodecimaldot]{babel}
\usepackage{lmodern}
\usepackage{geometry}
\usepackage{amsmath,amssymb}
\usepackage{booktabs}
\usepackage{longtable}
\usepackage{array}
\usepackage{xcolor}
\usepackage{hyperref}
\usepackage{enumitem}
\usepackage{microtype}
\usepackage{fancyhdr}
\geometry{margin=2.2cm}
\definecolor{darkblue}{HTML}{1E3A8A}
\hypersetup{colorlinks=true,linkcolor=darkblue,urlcolor=darkblue}
\pagestyle{fancy}
\fancyhf{}
\lhead{Investment Research AI}
\rhead{Metodología de evaluación}
\cfoot{\thepage}
\setlength{\parindent}{0pt}
\setlength{\parskip}{0.6em}

\title{\textbf{Metodología de Evaluación de Acciones}\\
\large Criterios, fórmulas, indicadores y ponderación del puntaje final}
\author{Investment Research AI}
\date{\today}

\begin{document}
\maketitle

\begin{abstract}
Este documento describe la metodología de evaluación cuantitativa utilizada por la plataforma Investment Research AI. Se detallan los 14 indicadores financieros actualmente empleados, sus fórmulas, rangos de evaluación, pesos y contribución al puntaje final. Los datos provienen principalmente de Yahoo Finance mediante \texttt{yfinance}. El score debe interpretarse como una herramienta de apoyo al análisis, no como una recomendación automática de compra o venta.
\end{abstract}

\tableofcontents
\newpage

\section{Fórmula general del puntaje}

Cada criterio disponible recibe un puntaje interno:
\[
S_i=
\begin{cases}
100, & \text{Cumple}\\
55, & \text{Revisar}\\
15, & \text{No cumple}
\end{cases}
\]

Los criterios sin dato no participan en el promedio. El score final es:
\[
\boxed{
Score=
\frac{\sum_{i=1}^{n} S_iW_i}
{\sum_{i=1}^{n}W_i}
}
\]
donde $S_i$ es el puntaje del criterio y $W_i$ su peso.

\subsection{Clasificación final}
\begin{center}
\begin{tabular}{ll}
\toprule
\textbf{Score} & \textbf{Clasificación}\\
\midrule
$\geq 80$ & Oportunidad prioritaria\\
$65$--$79.9$ & En seguimiento\\
$50$--$64.9$ & Neutral\\
$<50$ & Precaución\\
\bottomrule
\end{tabular}
\end{center}

\section{Pesos de los 14 criterios}

\begin{center}
\begin{tabular}{lr}
\toprule
\textbf{Criterio} & \textbf{Peso}\\
\midrule
Relación P/E & 10\%\\
Crecimiento de ingresos & 9\%\\
Crecimiento de ganancias & 9\%\\
Capitalización de mercado & 8\%\\
Potencial al precio objetivo & 8\%\\
Ventas totales & 7\%\\
ROE & 7\%\\
Margen operativo & 7\%\\
Deuda / patrimonio & 7\%\\
Flujo de caja libre & 7\%\\
Free float & 6\%\\
ROA & 5\%\\
Razón corriente & 5\%\\
Beta & 5\%\\
\midrule
\textbf{Total} & \textbf{100\%}\\
\bottomrule
\end{tabular}
\end{center}

\section{Indicadores, fórmulas y reglas}

\subsection{Capitalización de mercado --- 8\%}
\[
\boxed{Market\ Cap = Precio\ por\ acción \times Acciones\ en\ circulación}
\]
La plataforma expresa el valor en miles de millones:
\[
MarketCap_B=\frac{MarketCap}{10^9}
\]

\textbf{Regla:}
\begin{itemize}[leftmargin=*]
\item Cumple: $\geq 2$ mil millones USD.
\item Revisar: $0.3$--$1.999$ mil millones USD.
\item No cumple: $<0.3$ mil millones USD.
\end{itemize}

\textbf{Interpretación:} aproxima el tamaño, liquidez, estabilidad y cobertura institucional de la compañía.

\subsection{Relación P/E --- 10\%}
\[
\boxed{P/E=\frac{Precio\ por\ acción}{Ganancia\ por\ acción\ (EPS)}}
\]

La plataforma utiliza \texttt{trailingPE} y, si no está disponible, \texttt{forwardPE}.

\textbf{Regla:}
\begin{itemize}[leftmargin=*]
\item Cumple: $20\leq P/E\leq25$.
\item Revisar: $10\leq P/E\leq40$, excluyendo el rango ideal.
\item No cumple: $P/E<10$ o $P/E>40$.
\end{itemize}

\textbf{Interpretación:} mide cuánto paga el mercado por cada unidad de beneficio. Debe interpretarse junto con crecimiento y sector.

\subsection{Ventas totales --- 7\%}
\[
\boxed{Revenue_M=\frac{Revenue}{10^6}}
\]

\textbf{Regla:}
\begin{itemize}[leftmargin=*]
\item Cumple: $\geq1{,}000$ millones USD.
\item Revisar: $100$--$999.99$ millones USD.
\item No cumple: $<100$ millones USD.
\end{itemize}

\textbf{Interpretación:} ayuda a distinguir empresas con escala comercial de negocios todavía pequeños.

\subsection{Free Float --- 6\%}
\[
\boxed{
Free\ Float(\%)=
\frac{Float\ Shares}{Shares\ Outstanding}\times100
}
\]

\textbf{Regla:}
\begin{itemize}[leftmargin=*]
\item Cumple: $\geq40\%$.
\item Revisar: $20\%$--$39.99\%$.
\item No cumple: $<20\%$.
\end{itemize}

\textbf{Interpretación:} un float muy bajo puede implicar menor liquidez, mayor volatilidad y mayor concentración de propiedad.

\subsection{Precio objetivo promedio y potencial --- 8\%}

El precio objetivo mostrado corresponde al \texttt{targetMeanPrice} disponible en Yahoo Finance, es decir, el precio objetivo promedio de analistas.

\[
\boxed{
Upside(\%)=
\left(
\frac{Target\ Price-Current\ Price}{Current\ Price}
\right)\times100
}
\]

\textbf{Regla:}
\begin{itemize}[leftmargin=*]
\item Cumple: $\geq15\%$.
\item Revisar: $0\%$--$14.99\%$.
\item No cumple: $<0\%$.
\end{itemize}

\textbf{Interpretación:} incorpora la expectativa del consenso de analistas. No constituye una predicción propia de la plataforma ni garantiza rendimiento futuro.

\subsection{Crecimiento de ingresos --- 9\%}
\[
\boxed{
Revenue\ Growth(\%)=
\frac{Revenue_t-Revenue_{t-1}}{Revenue_{t-1}}\times100
}
\]

Yahoo Finance entrega el valor en forma decimal y la plataforma lo transforma:
\[
Growth_{\%}=Growth_{decimal}\times100
\]

\textbf{Regla:}
\begin{itemize}[leftmargin=*]
\item Cumple: $\geq10\%$.
\item Revisar: $0\%$--$9.99\%$.
\item No cumple: $<0\%$.
\end{itemize}

\textbf{Interpretación:} mide si el negocio está expandiendo sus ventas.

\subsection{Crecimiento de ganancias --- 9\%}
\[
\boxed{
Earnings\ Growth(\%)=
\frac{Earnings_t-Earnings_{t-1}}{Earnings_{t-1}}\times100
}
\]

\textbf{Regla:}
\begin{itemize}[leftmargin=*]
\item Cumple: $\geq10\%$.
\item Revisar: $0\%$--$9.99\%$.
\item No cumple: $<0\%$.
\end{itemize}

\textbf{Interpretación:} permite verificar si el crecimiento comercial se traduce también en crecimiento de beneficios.

\subsection{ROE --- 7\%}
\[
\boxed{
ROE=
\frac{Utilidad\ neta}{Patrimonio\ promedio}\times100
}
\]

\textbf{Regla:}
\begin{itemize}[leftmargin=*]
\item Cumple: $\geq15\%$.
\item Revisar: $8\%$--$14.99\%$.
\item No cumple: $<8\%$.
\end{itemize}

\textbf{Interpretación:} mide la rentabilidad generada sobre el capital de los accionistas.

\subsection{ROA --- 5\%}
\[
\boxed{
ROA=
\frac{Utilidad\ neta}{Activos\ totales\ promedio}\times100
}
\]

\textbf{Regla:}
\begin{itemize}[leftmargin=*]
\item Cumple: $\geq7\%$.
\item Revisar: $3\%$--$6.99\%$.
\item No cumple: $<3\%$.
\end{itemize}

\textbf{Interpretación:} mide la eficiencia con que la compañía utiliza sus activos.

\subsection{Margen operativo --- 7\%}
\[
\boxed{
Operating\ Margin=
\frac{Operating\ Income}{Revenue}\times100
}
\]

\textbf{Regla:}
\begin{itemize}[leftmargin=*]
\item Cumple: $\geq15\%$.
\item Revisar: $5\%$--$14.99\%$.
\item No cumple: $<5\%$.
\end{itemize}

\textbf{Interpretación:} muestra qué proporción de las ventas queda después de los costos y gastos operativos.

\subsection{Deuda / Patrimonio --- 7\%}
\[
\boxed{
D/E=\frac{Deuda\ total}{Patrimonio}
}
\]

\textbf{Regla:}
\begin{itemize}[leftmargin=*]
\item Cumple: $\leq100$.
\item Revisar: $100$--$200$.
\item No cumple: $>200$.
\end{itemize}

\textbf{Interpretación:} mide el apalancamiento financiero. Un valor muy alto puede implicar menor flexibilidad y mayor riesgo financiero.

\subsection{Razón corriente --- 5\%}
\[
\boxed{
Current\ Ratio=
\frac{Activos\ corrientes}{Pasivos\ corrientes}
}
\]

\textbf{Regla:}
\begin{itemize}[leftmargin=*]
\item Cumple: $\geq1.2$.
\item Revisar: $0.8$--$1.19$.
\item No cumple: $<0.8$.
\end{itemize}

\textbf{Interpretación:} mide la capacidad de cubrir obligaciones de corto plazo con activos corrientes.

\subsection{Flujo de caja libre --- 7\%}
\[
\boxed{
FCF=Operating\ Cash\ Flow-Capital\ Expenditures
}
\]

La plataforma lo expresa en millones:
\[
FCF_M=\frac{FCF}{10^6}
\]

\textbf{Regla:}
\begin{itemize}[leftmargin=*]
\item Cumple: $FCF\geq0$.
\item Revisar: $-1\leq FCF<0$ millones USD.
\item No cumple: $FCF<-1$ millón USD.
\end{itemize}

\textbf{Interpretación:} representa el efectivo disponible después de inversión operativa y puede utilizarse para reinversión, reducción de deuda, recompras o dividendos.

\subsection{Beta --- 5\%}
\[
\boxed{
\beta_i=
\frac{\operatorname{Cov}(R_i,R_m)}
{\operatorname{Var}(R_m)}
}
\]

\textbf{Regla:}
\begin{itemize}[leftmargin=*]
\item Cumple: $0.7\leq\beta\leq1.3$.
\item Revisar: $0\leq\beta\leq2$, excluyendo el rango ideal.
\item No cumple: $\beta<0$ o $\beta>2$.
\end{itemize}

\textbf{Interpretación:}
\begin{itemize}[leftmargin=*]
\item $\beta\approx1$: comportamiento similar al mercado.
\item $\beta>1$: mayor sensibilidad y volatilidad relativa.
\item $\beta<1$: menor sensibilidad relativa.
\end{itemize}

\section{Distribución conceptual del modelo}

\begin{center}
\begin{tabular}{lr}
\toprule
\textbf{Dimensión} & \textbf{Peso total}\\
\midrule
Crecimiento (ventas + earnings + revenue total) & 25\%\\
Valoración y expectativas (P/E + upside) & 18\%\\
Rentabilidad (ROE + ROA + margen operativo) & 19\%\\
Solvencia y caja (D/E + current ratio + FCF) & 19\%\\
Tamaño y liquidez (market cap + free float) & 14\%\\
Riesgo de mercado (beta) & 5\%\\
\bottomrule
\end{tabular}
\end{center}

\section{Ejemplo de cálculo del score}

Supóngase que una acción obtiene:

\begin{center}
\begin{tabular}{lrrrr}
\toprule
\textbf{Criterio} & \textbf{Peso} & \textbf{Estado} & \textbf{Score} & \textbf{Aporte}\\
\midrule
Capitalización & 8 & Cumple & 100 & 800\\
P/E & 10 & Revisar & 55 & 550\\
Ventas & 7 & Cumple & 100 & 700\\
Free float & 6 & Cumple & 100 & 600\\
Potencial & 8 & Cumple & 100 & 800\\
Revenue growth & 9 & Cumple & 100 & 900\\
Earnings growth & 9 & Revisar & 55 & 495\\
ROE & 7 & Cumple & 100 & 700\\
ROA & 5 & Revisar & 55 & 275\\
Margen operativo & 7 & Cumple & 100 & 700\\
D/E & 7 & Revisar & 55 & 385\\
Current ratio & 5 & Cumple & 100 & 500\\
FCF & 7 & Cumple & 100 & 700\\
Beta & 5 & Cumple & 100 & 500\\
\midrule
\textbf{Total} & \textbf{100} & & & \textbf{8,605}\\
\bottomrule
\end{tabular}
\end{center}

Por tanto:
\[
Score=\frac{8605}{100}=86.05
\]

La clasificación sería:
\[
\boxed{\text{Oportunidad prioritaria}}
\]

\section{Campos utilizados de Yahoo Finance}

La implementación actual utiliza principalmente:
\begin{itemize}[leftmargin=*]
\item \texttt{marketCap}
\item \texttt{trailingPE}
\item \texttt{forwardPE}
\item \texttt{totalRevenue}
\item \texttt{sharesOutstanding}
\item \texttt{floatShares}
\item \texttt{targetMeanPrice}
\item \texttt{revenueGrowth}
\item \texttt{earningsGrowth}
\item \texttt{returnOnEquity}
\item \texttt{returnOnAssets}
\item \texttt{operatingMargins}
\item \texttt{debtToEquity}
\item \texttt{currentRatio}
\item \texttt{freeCashflow}
\item \texttt{beta}
\end{itemize}

\section{Limitaciones}

\begin{enumerate}[leftmargin=*]
\item Los umbrales son reglas configuradas por la plataforma, no leyes universales de valoración.
\item P/E, márgenes, deuda y liquidez deben interpretarse según sector e industria.
\item El precio objetivo de analistas no garantiza rendimiento futuro.
\item Los criterios sin dato se excluyen del promedio ponderado.
\item El rango P/E de 20--25 refleja una preferencia específica de la metodología.
\item El FCF actualmente se evalúa sobre todo por ser positivo o negativo; podría mejorarse con FCF Yield.
\item Beta mide sensibilidad histórica al mercado, no riesgo total.
\end{enumerate}

\section{Conclusión}

El modelo actual sintetiza 14 indicadores de tamaño, liquidez, valoración, crecimiento, rentabilidad, solvencia, caja y riesgo. Su objetivo es priorizar empresas para análisis adicional y hacer transparente por qué una acción obtiene determinado score.

\end{document}
